\documentclass[11pt,a4paper]{article}
\usepackage[margin=0.85in,top=1in,bottom=1in]{geometry}
\usepackage{amsmath,amssymb,amsfonts}
\usepackage{graphicx}
\usepackage{float}
\usepackage{enumitem}
\usepackage{microtype}
\usepackage{caption}
\usepackage[hidelinks]{hyperref}
\usepackage[most,breakable]{tcolorbox}
\tcbuselibrary{breakable,skins}
\usepackage[utf8]{inputenc}
\usepackage[T1]{fontenc}
\usepackage{lmodern}
\captionsetup{font=small,labelfont=bf,skip=4pt}
\setlist{leftmargin=1.5em,topsep=2pt}

%% Breakable card — prevents content cutoff across pages
\newtcolorbox{solcard}[3]{
  breakable, enhanced,
  colback=white, colframe=gray!50, arc=2pt,
  title={\textbf{#1}\hfill\textsf{\small #3\,pts}},
  fonttitle=\normalsize\sffamily\bfseries,
  before upper={\small\textit{#2}\par\smallskip\hrule height 0.4pt\smallskip},
  top=4pt, bottom=6pt, left=7pt, right=7pt,
  parbox=false,
}

\title{\textbf{Pop-up bubbles --- Solution}}
\author{\normalsize X25 (2025) — English translation}
\date{}
\begin{document}\maketitle\vspace{-0.5em}
\begin{solcard}{A1}{What is the flow of the relative velocity vector of the fluid through an arbitrary closed surface that does not intersect the surface of the ball? Justify your answer. Note: the flow of the relative fluid velocity vector through a closed surface is the quantity $\Phi$, determined by the expression: \begin{equation*} \Phi=\oint_S\vec{v}_\text{rel}\bigl(\vec{r}\bigr)d\vec{S}{.} \end{equation*}}{0.30}

Note: the flow of the relative fluid velocity vector through a closed surface is the quantity $\Phi$, determined by the expression: \begin{equation*} \Phi=\oint_S\vec{v}_\text{rel}\bigl(\vec{r}\bigr)d\vec{S}{.} \end{equation*}


Due to the incompressibility of the liquid, its mass inside the selected surface is constant: $$0 = \frac{dm}{dt} = \displaystyle\oint \rho \vec{v}\cdot\vec{dS} = \rho \Phi$$

\par\smallskip\noindent\textbf{Answer:} $$\Phi = 0$$
\end{solcard}

\begin{solcard}{A2}{What is the normal component of the fluid velocity ${v}_{\text{rel},n}$ in the frame of reference associated with the ball?}{0.10}
\par\smallskip\noindent\textbf{Answer:} $v_{\text{rel}.n} = 0$
\end{solcard}

\begin{solcard}{A3}{What is the relative velocity of the fluid $\vec{v}_{\infty}$ at an infinite distance from the ball?}{0.10}
\par\smallskip\noindent\textbf{Answer:} The circulation theorem and the Gauss theorem (item A1) in differential form: $$\vec{\nabla} \times \vec{v}_{\text{rel}} = 0$$$$\vec{\nabla} \vec{v}_{\text{rel}} = 0$$Boundary condition: $v_{\text{rel},n} = 0$ $\vec{v}_{\text{rel}} (\infty) = -\vec{v_0}$
\end{solcard}

\begin{solcard}{A4}{Let us consider a superconducting body placed in a uniform magnetic field with induction $\vec{B}_0$. Write an expression for the circulation of the magnetic field $\vec{B}(\vec{r})$ along a contour that does not intersect the surface of the superconductor, and an expression for the flux of the magnetic field through a closed surface that does not intersect the surface of the superconductor. Also indicate the boundary conditions of the magnetic field at the boundary of the body and at an infinite distance from it.}{0.40}
\par\smallskip\noindent\textbf{Answer:} The circulation theorem and the Gauss theorem (item A1) in differential form: $$\vec{\nabla} \times  B = 0$$$$\vec{\nabla} \vec{B}  = 0$$Boundary condition: $B_{\text{rel},n} = 0$ $\vec{B}_{\text{rel}} (\infty) = \vec{B_0}$
\end{solcard}

\begin{solcard}{A5}{Let us place a superconducting ball of radius $R$ in a uniform magnetic field with induction $\vec{B_0}$. Find the magnetic field induction $\vec{B}(\vec{r})$ in the entire space outside the sphere. Express your answer using $\vec{B_0}, \vec{r}$ and $R$.}{0.80}

It is known that the field of a superconducting ball placed in a uniform magnetic field will coincide with the field of a magnetic dipole. The expression for the total magnetic field takes the form $$\vec{B}=\vec{B_0}+\frac{\mu_0}{4\pi}\left(\frac{3(\vec{m}\vec{r})\vec{r}}{r^5}-\frac{\vec{m}}{r^3}\right) $$where $m$ - magnetic moment. From the condition zero normal constitute magnetic field on surface sphere obtain $$(\vec{B}\vec{r})=(\vec{B_0}\vec{r})+\frac{\mu_0}{2\pi}\frac{(\vec{m}\vec{r})}{R^4}=0 $$From last relation find magnetic moment sphere $$\vec{m}=-\frac{2\pi R^3}{\mu_0}\vec{B_0} $$After finding the magnetic moment, the expression for the magnetic field takes the form:

\par\smallskip\noindent\textbf{Answer:} $$\vec{B}=\vec{B_0}-\frac{3R^3(\vec{B_0}\vec{r})\vec{r}}{2r^5}+\frac{\vec{B_0}R^3}{2r^3} $$
\end{solcard}

\begin{solcard}{A6}{Using the results of the previous paragraphs, find the dependence $\vec{v}_{\text{rel}}(\vec{r})$.}{0.20}

The system of equations on $\vec{v}_{\text{rel}}(\vec{r})$ coincides with the system of equations on $\vec{B}(\vec{r})$ up to the change $\vec{B}_0 = -\vec{v}_0$.

\par\smallskip\noindent\textbf{Answer:} $$\vec{v}_{\text{rel}}(\vec{r}) = -\vec{v_0}+\frac{3R^3(\vec{v_0}\vec{r})\vec{r}}{2r^5}-\frac{\vec{v_0}R^3}{2r^3} $$
\end{solcard}

\begin{solcard}{A7}{Let us consider the stationary motion of an incompressible fluid with zero viscosity through a thin tube. Let $p$ be the pressure in the liquid, $\varphi$ the potential energy per unit mass of the liquid, and $v$ the flow velocity in the pipe. Using the laws of conservation of energy and mass, show that for all points of the tube the relation (Bernoulli equation) holds: $$\varphi  + \frac{p}{\rho} + \frac{ v_\text{rel}^2}{2} = const$$}{0.40}
\par\smallskip\noindent\textbf{Answer:} In the reference frame of the ball, the distribution of velocities in space, and therefore the streamlines, do not depend on time. Let us select one streamline through which the mass flow rate $\frac{dm}{dt} = \rho \frac{dv}{dt}$ flows. Let us select two points (1) and (2) in this streamline. Consider the fluid contained between these two points. The power of the external pressure force at point (1) is equal to $P_1 = p_1 \frac{dV}{dt} = p_1  \frac{dm}{dt} / \rho $, and at point (2) is equal to $P_2 = - p_2 \frac{dV}{dt} = p_2  \frac{dm}{dt} / \rho $. In time $dt$, the work of pressure forces is spent on accelerating a piece of liquid with mass $dm$ from speed $v_1$ to speed $v_2$: $$(P_1 + P_2)\cdot dt = dm (\frac{v_2^2}{2} - \frac{v_1^2}{2})$$ $$p_1/ \rho - p_2 / \rho = \frac{v_2^2}{2} - \frac{v_1^2}{2}$$$$p_1 + \frac{\rho v_1^2}{2} = p_2 + \frac{\rho v_2^2}{2}$$Which was what needed to be proven.
\end{solcard}

\begin{solcard}{A8}{Assuming that the pressure at infinite distance from the ball is equal to $p_0$, find the dependence $p(\theta)$ on the surface of the ball. Express your answer using $p_0$, $\rho$, $v_0$, $\theta$.}{0.60}

Let us write down the Bernoulli equation obtained in paragraph A7 at two points - at infinity and on the surface of the ball: $$p(\theta) + \frac{\rho v^2_{\text{surf}}(\theta)}{2} =p_{0} + \frac{\rho v^2_{0}}{2} $$ $$p(\theta)  =p_{0} + \frac{\rho ( v^2_{0} - v^2_{\text{surf}}(\theta))}{2} $$in formula from part A6 substitute $r = R$: $$\vec{v}_{\text{surf}}(\vec{e}_r) = -\vec{v_0}+\frac{3(\vec{v_0}\vec{e}_r)\vec{e}_r}{2}-\frac{\vec{v_0}}{2} = \frac{3}{2} \big( -\vec{v_0}+(\vec{v_0}\vec{e}_r)\vec{e}_r\big)$$$$\vec{v}^2_{\text{surf}}(\vec{e}_r) = \frac{9}{4}\big(v_0^2 - 2(\vec{v_0}\vec{e}_r)^2 + (\vec{v_0}\vec{e}_r)^2 \big) = \frac{9}{4}\big(v_0^2 - (\vec{v_0}\vec{e}_r)^2 \big) = \frac{9}{4}v_0^2 \sin^2{\theta} $$

\par\smallskip\noindent\textbf{Answer:} $$p(\theta) = p_0 + \frac{\rho v_0^2}{2} (1 - \frac{9}{4}\sin^2{\theta})$$
\end{solcard}

\begin{solcard}{A9}{Calculate the resultant of the pressure forces acting on the surface of the ball. Express your answer in terms of $\vec{v_0}$, $R$, $\rho$.}{0.40}
\par\smallskip\noindent\textbf{Answer:} The dependence $p(\theta)$ is an even function, which means the force will not change if we start counting $\theta$ from the opposite pole of the ball. Therefore, there is no force component codirectional with $\vec{v}_0$. Other components are missing due to symmetry.
\end{solcard}

\begin{solcard}{A10}{Find the distribution of fluid velocity $\vec{v}(\vec{r})$ in the laboratory frame of reference.}{0.10}
\par\smallskip\noindent\textbf{Answer:} $$\vec{v}(\vec{r}) = \vec{v}_0 + \vec{v}_{\text{rel}}(\vec{r}) = \frac{3R^3(\vec{v_0}\vec{r})\vec{r}}{2r^5}-\frac{\vec{v_0}R^3}{2r^3}$$
\end{solcard}

\begin{solcard}{A11}{What is the kinetic energy of the liquid $E_k$ in the laboratory frame of reference? Express your answer in terms of $\rho$, $R$, $v_0$.}{0.80}

To calculate the kinetic energy of a liquid, it is necessary to calculate the integral $\int\frac{\rho v^2 dV}{2}$. For a given angle $\theta$, select a segment with angle $d\theta$ and thickness $dr$. The volume of this element is $2\pi r^2sin(\theta)~dr~d{\theta}$. Then $$W_to=\int\limits_R^{\infty} \int\limits_0^\pi {\rho\pi}r^2 \sin (\theta)drd{\theta}{v^2({\theta})} $$Find ${v^2({\theta})}$ from theorem cosine $${v^2({\theta})}=\frac{v^2\left(1+3\cos^2(\theta)\right)R^6}{4r^6} $$Combining the last expression with the already obtained integral give $$\int\limits_R^{\infty} \frac{dr}{r^4}=\frac{1}{3R^3} $$$$\int\limits_0^{\pi} \sin(\theta)(1+3\cos^2(\theta))d{\theta}=\int\limits_1^{-1} -(d\cos(\theta)+3\cos^2(\theta)d\cos(\theta))=4 $$

\par\smallskip\noindent\textbf{Answer:} $$W_to=\frac{\pi \rho R^3v^2}{3} $$
\end{solcard}

\begin{solcard}{A12}{Find the added mass $\mu$. Express your answer in terms of $\rho$, $R$.}{0.20}
\par\smallskip\noindent\textbf{Answer:} $$\mu=\frac{2W_to}{v^2}=\frac{2\pi \rho R^3}{3} $$
\end{solcard}

\begin{solcard}{A13}{Consider a ball moving with speed $\vec{v}_0$ and acceleration $\vec{a}_0$. Using the expression for the momentum $\vec{p}$ of the fluid, obtain an expression for the force $\vec{F}$ acting on the ball from the fluid. Express your answer using $\rho$, $R$ and $\vec{a}_0$.}{0.30}

According to Newton's first law, the force $-\vec{F}$ acts on water from the side of the ball. Newton's second law for liquid: $$-\vec{F} = \frac{d}{dt}(\vec{p}) = \mu \vec{a}_0$$.

\par\smallskip\noindent\textbf{Answer:} $$\vec{F} = -\mu \vec{a}_0$$
\end{solcard}

\begin{solcard}{A14}{Obtain the dependence of the curvature $K$ of the bubble surface on the angle $\theta$. Use $p_{in}$, $p_0$, $\rho$, $\theta$ and $v_0$. Note: the curvature $K$ of a surface is the sum of the inverse principal radii of curvature: $$ K=\frac{1}{R_1}+\frac{1}{R_2}$$ It is related to the Laplace pressure as: $$\Delta p = \sigma K$$}{0.30}

Note: the curvature $K$ of a surface is the sum of the inverse principal radii of curvature: $$ K=\frac{1}{R_1}+\frac{1}{R_2}$$


It is related to Laplace pressure as:


$$\Delta p = \sigma K$$


$$\Delta P = p_{in} - p_0 - \frac{\rho v_0^2}{2} + \frac{9\rho v_0^2}{8} \sin^2(\theta) = \sigma K$$

\par\smallskip\noindent\textbf{Answer:} $$K= \frac{1}{\sigma} \big( p_{in} - p_0 - \frac{\rho v_0^2}{2} + \frac{9\rho v_0^2}{8} \sin^2(\theta) \big)$$
\end{solcard}

\begin{solcard}{B1}{Taking into account the smallness of $\varepsilon$, the values $R_1$ and $R_2$ can be represented in the form: \begin{equation*} R_1(\theta)\approx R(1+A_1\varepsilon-B_1\varepsilon\sin^2\theta)\qquad R_2\approx R(1+A_2\varepsilon-B_2\varepsilon\sin^2\theta){.} \end{equation*} Find $B_1$ and $B_2$. Further, the terms $A_1\varepsilon$ and $A_2\varepsilon$ can be neglected. If you were unable to solve this point, then you can further assume that $B_1=B_2=1$. This does not lead to the loss of subsequent points.}{0.60}

Let's find the radius of curvature in the $(x, y)$ plane. Formula for radius of curvature in 2D coordinates: $$R = \frac{(1 + (dy/dx)^2)^{3/2}}{d^2y/dx^2}$$ $$\frac{dy}{dx} = \tan{\theta}$$Derivative express through $\theta$: $$\frac{d^2y}{dx^2} = \frac{1}{R \cos^3{\theta}}$$ Then we stretched the same body by a factor of $k$ along the $x$ axis. First and second derivative $\frac{dy}{d x}$ and $\frac{d^2y}{dx^2}$ at each point will change by a factor of $k$. $$R_1(\theta) = \frac{(1 + (dy'/dx)^2)^{3/2}}{d^2y'/dx^2} = \frac{(1 + (kdy'/dx)^2)^{3/2}}{kd^2y/(dx)^2} = \frac{(1 + k^2 \cdot (dy/dx)^2)^{3/2}}{k\cdot d^2y/dx^2} = $$ $$= k R \cos^3{\theta}\big(1 + k^2 \cdot \tan{\theta}^2\big)^{3/2} = \frac{R}{k} \cdot (\cos^2{\theta} + k^2 \sin^2{\theta})^\frac{3}{2}$$


Counting $R \approx a$: $$R_1(\theta)= \frac{a}{k} \cdot (\cos^2{\theta} + k^2 \sin^2{\theta})^\frac{3}{2} = \frac{a}{1-\varepsilon} \cdot (\cos^2{\theta} + (1 - \varepsilon)^2 \sin{\theta}^2)^\frac{3}{2} \approx$$ $$\approx a (1 + \varepsilon) (\cos^2{\theta} + (1 - 2\varepsilon) \sin^2{\theta})^\frac{3}{2} = a (1 + \varepsilon) (1 - 2\varepsilon  \sin^2{\theta})^\frac{3}{2} \approx$$$$\approx a (1 + \varepsilon) (1 - 3\varepsilon  \sin^2{\theta}) \approx a(1 + \varepsilon - 3\varepsilon\sin^2{\theta}) \approx R (1 + \varepsilon - 3\varepsilon\sin^2{\theta}) $$Obtain: $A_1 = 1$; $B_1 = 3$;


Now let us find the radius of curvature in the plane $\alpha$, perpendicular to $x, y$ and containing the normal to the surface at point $A'$. Then we can use Laplace's formula for surface tension pressure.


In the figure above, $\alpha$ and $\beta$ are planes perpendicular to the plane of the figure. At point $A'$, the curve obtained by cutting the ellipsoid by plane $\beta$ has a center of curvature $O_1$. Similarly, the curve obtained by cutting an ellipsoid by plane $\alpha$ has a center of curvature $O_2$. Meunier's theorem states that since $O_2$ is the principal center of curvature, then its projection onto the $\beta$ plane is the center of curvature in that plane - that is, $O_1$. $$R_{2} = R_{\alpha} = R_{\beta} / \sin{\theta'}$$$R_{\beta} = R \sin{\theta}$ - radius of curvature in the $\beta$ plane. It does not change when stretched.  Then: $$R_2 = R \cdot \cfrac{\sin{\theta}}{\sin{\theta'}}$$At the same time $\theta'$ will also be the angle of inclination of the tangent to the ellipse at the point $A'$: $$\tan{\theta'} = (|\frac{dy}{dx}|)_{A'} =  k \cdot (|\frac{dy}{dx}|)_{A} = k \tan{\theta}$$Then: $$R_2(\theta) = \frac{R\sin{\theta}}{\sin{\theta'}} = R \frac{\tan{\theta}}{\sqrt{1 + \tan^2{\theta}}} \cdot \frac{\sqrt{1 + \tan^2{\theta'}}}{\tan{\theta'}} =$$ $$= R \frac{\tan{\theta}}{\sqrt{1 + \tan^2{\theta}}} \cdot \frac{\sqrt{1 + k^2\tan^2{\theta}}}{k\tan{\theta}} = \frac{R}{k} \sqrt{\frac{1 + k^2\tan^2{\theta}}{1 + \tan^2{\theta}}}  = \frac{R}{k} \sqrt{\cos^2{\theta} + k^2 \sin^2{\theta}}$$


Counting $R \approx a$: $$R_2(\theta) = \frac{a}{k} \sqrt{\cos^2{\theta} + k^2 \sin^2{\theta}} = \frac{a}{1 - \varepsilon} \sqrt{\cos^2{\theta} + (1-\varepsilon)^2 \sin^2{\theta}} \approx$$$$\approx a(1 + \varepsilon)\sqrt{\cos^2{\theta} + (1 - 2\varepsilon) \sin^2{\theta}} = a(1 + \varepsilon)\sqrt{1 - 2\varepsilon \sin^2{\theta}} \approx$$$$\approx a(1+\varepsilon)(1 - \varepsilon\sin^2{\theta}) \approx a(1 + \varepsilon - \varepsilon \sin^2{\theta}) \approx R(1 + \varepsilon - \varepsilon \sin^2{\theta})$$Obtain: $A_2 = 1$; $B_2 = 1$;

\begin{figure}[H]
  \centering
  \includegraphics[width=0.72\textwidth]{T4_S_EN_assets/fig1.jpg}
\end{figure}
\end{solcard}

\begin{solcard}{B2}{Find the dependence of the curvature of the ellipsoid $K(\theta)$ in first order in $\varepsilon$. Express your answer using $R$, $\theta$.}{0.20}

$$K = \frac{1}{R(1 - B_1 \varepsilon \sin^2{\theta})} + \frac{1}{R(1 - B_2 \varepsilon \sin^2{\theta})} \approx \frac{1}{R}\big(1 +  B_1 \varepsilon \sin^2{\theta}+1+  B_2 \varepsilon \sin^2{\theta} \big) =$$ $$= \frac{1}{R}\big(2 + ( B_1 +B_2) \varepsilon \sin^2{\theta}  \big) = \frac{1}{R}\big(2 + 4\varepsilon \sin^2{\theta}  \big)$$

\end{solcard}

\begin{solcard}{B3}{What is the value of $\varepsilon$ for the bubble discussed in A14? Express your answer in terms of bubble speed $v_0$, bubble radius $R$, surface tension coefficient $\sigma$ and liquid density $\rho$.}{0.60}

Let's use the formula for Laplace Pressure: $$\Delta p = \sigma K$$According to part A13: $$\Delta P = p_{in} - p_0 - \frac{\rho v_0^2}{2} + \frac{9\rho v_0^2}{8} \sin^2(\theta) = \sigma /R \cdot ( 2 + 4\varepsilon \sin^2{\theta})$$Since equality hold for all angle $\theta$, that coefficient before term $\sin^2{\theta}$ must is equal to: $$\frac{9\rho v_0^2}{8} = \frac{4\varepsilon \sigma}{R}$$

\par\smallskip\noindent\textbf{Answer:} $$ \varepsilon   = \frac{9R\rho v_0^2}{32\sigma}$$
\end{solcard}

\begin{solcard}{B4}{What is the eccentricity $e$ of the bubble discussed in A14? Express your answer in terms of bubble speed $v_0$, bubble radius $R$, surface tension coefficient $\sigma$ and liquid density $\rho$.}{0.30}

$$e = \frac{\sqrt{a^2 - b^2}}{a} = \sqrt{1-(1-\varepsilon)^2} \approx \sqrt{2\varepsilon} =  \sqrt{\frac{9R\rho v_0^2}{16\sigma}}$$

\end{solcard}

\begin{solcard}{B5}{A spherical bubble with an initial diameter $d = 2\text{mm}$ rises in water with density $\rho = 10^3 \text{kg}/\text{s}$ and surface tension $\sigma = 72 \cdot 10^{-3} ~\text{N}/\text{m}$. After taking the necessary measurements, calculate the speed of the bubble in the last three photographs. Also calculate the acceleration of the bubble $a$.}{0.80}

$$v_2^2=v_3^2-2al_{23}$$ $l_{23}$ can be obtained by comparing the distance with the horizontal diameter of the bubble on the first frames. $$l_{23}=2.8 mm$$

\par\smallskip\noindent\textbf{Answer:} $$v_1=17\text{cm}/\text{c}, v_2=26\text{cm}/\text{c}, v_3=32\text{cm}/\text{c}, a = 6 \text{m}/\text{c}^2$$
\end{solcard}

\begin{solcard}{C1}{What is its acceleration $a$ at the initial moment of time? Express your answer using $g$.}{0.60}

With neglecting air mass: $$0 = 0 + \frac{4}{3}\pi R^3 \rho g - \mu a - \pi R^2 \rho v^2$$in initial moment time velocity bubble equal zero, and means: $$a = \frac{\frac{4}{3}\pi R^3 \rho}{\mu} g = 2g$$

\par\smallskip\noindent\textbf{Answer:} $$a  = 2g$$
\end{solcard}

\begin{solcard}{C2}{Find the steady speed of the bubble $v_{\max}$. What is the value of $\varepsilon$ at steady speed? Express your answers using $R$, $g$, $\rho$ and $\sigma$.}{0.30}

The bubble's acceleration is zero: $$0 = 0 + \frac{4}{3}\pi R^3 \rho g - \pi R^2 \rho v_{\max}^2$$$$v_{\max} = \sqrt{\frac{4}{3}Rg}$$$$ \varepsilon   = \frac{9R\rho v_{\max}^2}{32\sigma} = \frac{9R\rho }{32\sigma} \cdot \frac{4}{3}Rg = \frac{3R^2 \rho g}{8\sigma}$$

\end{solcard}

\begin{solcard}{C3}{Calculate the numerical values of $v_{\max}$ and $\varepsilon$ for $\rho = 10^3~ {\text{kg}}/{\text{m}^3}$, $\sigma = 72 \cdot 10^{-3} ~\text{N}/\text{m}$, $g =10 ~\text{m}/ \text{s}^{2}$ and $R = 1~\text{mm}$. Is the selected model applicable to the specified radius?}{0.60}
\par\smallskip\noindent\textbf{Answer:} $$v_{max} = \sqrt{\frac{4}{3}Rg} = 12 cm/\mathrm{s}$$$$\varepsilon = \dfrac{3R^2\rho g}{8\sigma} = 5.2 \cdot 10^{-2} \ll 1$$$$\varepsilon \ll 1 \Rightarrow \text{Model is applicable}$$
\end{solcard}

\begin{solcard}{D1}{Prove that the water pressure at points on the surface of the bubble is constant.}{0.20}
\par\smallskip\noindent\textbf{Answer:} The pressure above the bubble is calculated using the formula $$p  = p_{in} - \sigma K$$ where both terms are constant.
\end{solcard}

\begin{solcard}{D2}{Find the pressure relationship above the surface of bubble $p(\theta)$ using $\rho$, $v_0$, $g$, $R$, $\theta$ and $h$.}{1.00}

Let us write the Bernoulli equation in the bubble reference frame for the streamline from the surface of the reservoir to a certain point on the bubble: $$p_0 + \rho g h + \frac{\rho v_0^2}{2} = p(\theta) + \frac{\rho v^2(\theta)}{2} + \rho g R(\cos{\theta} - 1)$$We this can do, since the linear current in this reference frame is constant due to $h \gg R$ (we neglect the edge effects that arise on the surface of the reservoir) According to the condition, the velocity distribution is the same as in part A:

\par\smallskip\noindent\textbf{Answer:} $$p(\theta) = p_0 +\rho g h + \rho \frac{v_0^2}{2} \rho g R(1 - \cos{\theta}) -  \frac{9}{8}\rho v_0^2 \sin^2{\theta}$$
\end{solcard}

\begin{solcard}{D3}{Derive an expression for the steady-state speed of the bubble $v_0$. Calculate the value for $R = 30~\text{cm}$ and $g = 10~\text{m}/\text{s}^2$.}{0.80}

Let us expand the formula from the previous paragraph to the second order in angle $\theta$: $$p(\theta) = p_0 +\rho g h + \rho \frac{v_0^2}{2} + \rho g R \cdot \theta^2 /2 -  \frac{9}{8}\rho v_0^2 \cdot \theta^2$$Using the fact of constancy of the pressure, we equate the coefficients before $\theta^2$: $$  \rho g R  / 2= \frac{9}{8}\rho v_0^2 $$Answer:

\par\smallskip\noindent\textbf{Answer:} $$v_0 = \sqrt{\dfrac{4}{9}Rg}=1.2 m/c$$
\end{solcard}

\end{document}